% Nouveau EGA V (cote 25): archive cote 25, site page 13 (PDF page 14).
% typescript. Draft compared: none.
% AI-transcribed by Claude Opus 5.5 (claude-opus-5-5, adaptive thinking, effort high), Message Batches API, batch ega25b (msgbatch_01A7n5YuXubUPNK1zsepi8RG), 2026-11-13; not yet checked by eye.
\scanpage{25/13}
\mspage{\hand{10}}

sont linéairement indépendantes. Or ces différentielles ne sont autres
que les sections de $(\underline{\hand{\Omega}}^1_{U/k})\otimes_k k(\ )$ suivantes:
\[ df_1, \ldots, df_p , \hand{\sum_i} S_i dT_i \ . \]
Comme les $df_i$ sont linéairement indépendants en tout points de U, et
que les $dT_i$ forment une base \var{\inserted[typed]{de $\underline{\hand{\Omega}}^1_{U/S}$}}{de $\underline{\hand{\Omega}}^1_{U/S}$}{insertion}{ega25_25_013-v0001} en tout point de U, \var{\xout{\illegible{1}}}{}{typewriter strike-out}{ega25_25_013-v0002} et à forti-%
ori un système de générateurs, on en conclut immédiatement l'indépendan-%
ce linéaire des quantités écrites en tout point de $U_{k(\hand{\eta})}$ , du moins
lorsque $p\leq n$, i.e. si \var{\inserted[typed]{$\underline{E} =$}}{$\underline{E} =$}{insertion}{ega25_25_013-v0003} $\underline{\hand{\Omega}}^1_{U/k} / \hand{\sum} \underline{O}_U df_i \neq 0$ : c'est un petit
lemme sur une famille de \unsure{gébérateurs|générateurs} \var{\inserted[typed]{$\hand{\unsure{a_i}}$ $(1 \hand{\leq} i \hand{\leq} n)$}}{$\hand{\unsure{a_i}}$ $(1 \hand{\leq} i \hand{\leq} n)$}{insertion}{ega25_25_013-v0004} d'un Module loc libre $\underline{E}$ non
nul, alors $\hand{\sum} S_i a_i$ considéré comme section de $\underline{E}\otimes_k k(\hand{\eta})$ ne s'annule
en aucun point. D(autre part, le cas où p=n est trivial, car alors
$Y_{\hand{\eta}} = \emptyset$ .

\uline{Corollaire} 2.13. Soit Z l'ensemble des points de X en lesquels X
n'est pas lisse sur k. Pour que Y soit lisse sur $k(\hand{\eta})$, il f et s
que Z soit fini. En particulier, si X est lisse, il en est de même de
Y .

Résulte de 2.12. et 2.3. Plus généralement, on obtient:

\var{\margin{\unsure{un ramifié|non ramifié} \uline{ou} k de car 0}}{\placedmargin{\unsure{un ramifié|non ramifié} \uline{ou} k de car 0}}{marginal note placed}{ega25_25_013-v0005}
\uline{\unsure{Théoèème|Théorème}} 2.1\unsure{\var{\over{3}{4}}{4}{writing over}{ega25_25_013-v0006}|\var{\xout{3}}{}{typewriter strike-out}{ega25_25_013-v0007}4}. Soi\var{\xout{e}}{}{typewriter strike-out}{ega25_25_013-v0008}nt \unsure{y} un point de $\unsure{\var{\over{X}{Y}}{Y}{writing over}{ega25_25_013-v0009}|\var{\xout{X}}{}{typewriter strike-out}{ega25_25_013-v0010}Y}_{\hand{\eta}}$ , x son image dans X. Soit
\var{\inserted[typed]{$P(\underline{A},K)$}}{$P(\underline{A},K)$}{insertion}{ega25_25_013-v0011} \var{\deleted{P}}{}{deletion}{ega25_25_013-v0012} l'une quelconque des propriétés suivantes pour une\var{\xout{\illegible{1}}}{}{typewriter strike-out}{ega25_25_013-v0013} algèbre
locale noethérienne sur un corps K:

(i) A est géométriquement régulier sur K.

(ii) A est géométriquement $(R_k)$ sur K.

(iii) A est séparable sur K.

(iv) A est \var{\inserted[typed]{géométriquement}}{géométriquement}{insertion}{ega25_25_013-v0014} normal sur K.

\var{\inserted{Alors $P(\underline{O}_{X,x}, k) \Longleftrightarrow P(\underline{O}_{Y_\eta,y}, k(\eta))$}}{Alors $P(\underline{O}_{X,x}, k) \Longleftrightarrow P(\underline{O}_{Y_\eta,y}, k(\eta))$}{insertion}{ega25_25_013-v0015}
\var{\xout{(x)xxxxxAxxxx}}{}{typewriter strike-out}{ega25_25_013-v0016} En effet, compte tenu du par.6 , (iii) et (iv)
résultent de (ii) conjugué avec 2.8. (ii). D'autre part, (i)
a déjà été prouvé dans 2.12. , et il reste à ramener (ii) à (i). Pour

\mspage{\hand{13}}
